% =============================================================================
\section{Trajectory Generation from Waypoints: Mission Manager}
\label{sec:mission}
% =============================================================================

\subsection{GPS $\to$ local ENU conversion}
\label{sec:gps2enu}

All guidance math (Section~\ref{sec:ilos}) operates in a local, metric East-North-Up (ENU) frame, while missions are authored in geodetic coordinates (latitude/longitude). \texttt{MissionManager} converts every waypoint once, on receipt of the local datum $(\varphi_0,\lambda_0)$ (Section~\ref{sec:ekf}), using the same flat-Earth (equirectangular) approximation used throughout the stack (accurate for ranges $\lesssim 10$\,km):
\begin{align}
x &= (\lambda - \lambda_0)\,\cos(\varphi_0)\cdot k, \\
y &= (\varphi - \varphi_0)\cdot k, \qquad k = 111319.5\ \text{m/deg (ArduPilot constant} \approx R_\oplus\cdot\pi/180).
\end{align}
This gives \texttt{MissionManager} a sequence of local ENU waypoints, which is what is actually handed to the ILOS guidance law leg-by-leg (Section~\ref{sec:ilos}).

\subsection{Mission finite-state machine}

\texttt{MissionManager} sequences a deploy/retrieve-style mission as a 6-state FSM, shown in Figure~\ref{fig:missionfsm}, that advances from waypoint to waypoint and reacts to a small set of service calls and arrival events.

\begin{figure}[H]
\centering
\begin{tikzpicture}[>=Stealth,node distance=14mm,every node/.style={font=\scriptsize}]
\tikzstyle{st}=[draw,rounded corners,fill=green!7,minimum width=3.0cm,minimum height=0.9cm,align=center]
\node[st] (idle) {IDLE\_AT\_BASE};
\node[st,right=32mm of idle] (deploy) {TRANSITING\\TO\_DEPLOY};
\node[st,below=of deploy] (atdeploy) {AT\_DEPLOY\\POSITION};
\node[st,below=of idle] (return) {RETURN\_TO\_BASE};
\node[st,right=32mm of atdeploy] (landtrans) {TRANSITING\_TO\\LANDING\_POS};
\node[st,below=of landtrans] (landing) {UAV\_LANDING};

\draw[->] (idle) -- node[above]{svc: prepare\_takeoff} (deploy);
\draw[->] (deploy) -- node[right]{waypoints reached} (atdeploy);
\draw[->] (atdeploy) -- node[above]{svc: tookoff} (landtrans);
\draw[->] (landtrans) -- node[right]{landing pos.\ reached} (landing);
\draw[->] (landing) -- node[above]{svc: landed} (return);
\draw[->] (return) -- node[left]{waypoints reached} (idle);
\draw[->,red,dashed] (landtrans) -- node[right,red]{svc: emergency\_landing} (landing);
\end{tikzpicture}
\caption{MissionManager finite-state machine: waypoints are consumed sequentially within each \texttt{TRANSITING\_*} leg, and the FSM advances once the leg's waypoints are reached or the corresponding service is called.}
\label{fig:missionfsm}
\end{figure}

\todored{Confirm/describe the cooperating UAV platform and CONOPS behind the deploy/retrieve service names (\texttt{prepare\_takeoff}, \texttt{tookoff}, \texttt{landed}, \texttt{emergency\_landing}), which indicate this mission profile was built for a cooperative USV$\to$UAV deploy/recovery operation.}

\subsection{Example missions}

Table~\ref{tab:missions} lists the two GPS mission definitions currently checked in, each specified as an ordered sequence of waypoints that \texttt{MissionManager} converts to ENU (Section~\ref{sec:gps2enu}) and feeds leg-by-leg to the ILOS guidance law.

\begin{table}[H]
\centering
\caption{Mission waypoint sets (WGS84 degrees).}
\label{tab:missions}
\begin{tabular}{@{}lll@{}}
\toprule
Mission & Leg & Waypoints (lat, lon) \\
\midrule
\multirow{3}{*}{Square (FEUP)} & Deploy & (41.685117, $-8.838916$) $\to$ (41.685399, $-8.839086$) $\to$ (41.685684, $-8.838855$)\\
 & Landing & (41.685399, $-8.839086$) \\
 & Return & (41.685117, $-8.838916$) \\
\midrule
\multirow{4}{*}{Viana buoys} & Deploy & (41.684793,$-8.839555$)$\to$(41.684922,$-8.839550$)$\to$(41.685022,$-8.839395$)\\
 & Deploy pos.\ & (41.685022, $-8.839395$)\\
 & Landing & (41.686183, $-8.839467$) \\
 & Return & reverse of deploy leg \\
\bottomrule
\end{tabular}
\end{table}
